\documentclass[conference]{IEEEtran}
\IEEEoverridecommandlockouts
\usepackage{cite}
\usepackage{amsmath,amssymb,amsfonts}
\usepackage{algorithmic}
\usepackage{graphicx}
\usepackage{textcomp}
\usepackage{xcolor}
\def\BibTeX{{\rm B\kern-.05em{\sc i\kern-.025em b}\kern-.08em
    T\kern-.1667em\lower.7ex\hbox{E}\kern-.125emX}}
\begin{document}

\title{An Educational Framework for AI-Driven RF Signal Processing on FPGA}

\author{\IEEEauthorblockN{Jude Eschete}
\IEEEauthorblockA{\textit{Dept. of Electrical and Computer Engineering} \\
\textit{Stevens Institute of Technology}\\
Hoboken, NJ, USA \\
JEschete@stevens.edu}
\and
\IEEEauthorblockN{Project Advisor: Bernard Yett}
\IEEEauthorblockA{\textit{Dept. of Electrical and Computer Engineering} \\
\textit{Stevens Institute of Technology}\\
Hoboken, NJ, USA \\
byett1@stevens.edu}
}

\maketitle

\begin{abstract}
University electrical engineering programs typically deliver instruction in FPGA digital design, RF communications, embedded systems, and machine learning as independent courses, leaving graduates without structured experience integrating these disciplines. This paper presents an educational framework built around a benchtop hardware platform in which STM32 microcontrollers paired with RFM95W LoRa transceivers transmit RF signals that are received and processed in real time by a Digilent Nexys A7-100T FPGA, with collected signal data feeding a machine learning classification pipeline. A five-module progressive curriculum sequences students from hardware bring-up and embedded firmware development through FPGA RTL design, system integration, and AI-driven signal analysis. Each module makes the interfaces between disciplines explicit teaching objectives, enabling students to reason about how design decisions in one domain propagate constraints into adjacent ones.
\end{abstract}

\begin{IEEEkeywords}
FPGA, LoRa, machine learning, RF signal processing, embedded systems, curriculum design, signal classification
\end{IEEEkeywords}

\section{Introduction}

Modern electrical and computer engineering practice increasingly demands fluency across multiple technical domains simultaneously. Defense, telecommunications, and aerospace systems routinely integrate real-time RF data acquisition, hardware-accelerated digital processing, embedded firmware, and intelligent signal analysis within a single pipeline. University curricula, however, frequently deliver these disciplines in isolation. A student may learn HDL synthesis and timing closure in a digital design course, modulation theory and channel models in a communications course, interrupt-driven firmware in an embedded systems course, and classification on benchmark datasets in a machine learning course—without encountering a structured experience that requires synthesizing these skills into a working system.

This isolation creates a gap between academic preparation and professional practice. Engineers who understand each domain individually often struggle with the interface problems that arise when combining them: how RF modulation parameters affect FPGA buffer sizing, how reconfigurable logic resource constraints shape a data collection pipeline for machine learning, or how real-world RF propagation deviates from classroom channel models. These cross-disciplinary concerns are not incidental details but core engineering competencies that no single-domain course addresses.

Prior research has contributed to each component of this problem in isolation. Educational frameworks for FPGA-based digital design have shown that cloud-based, open-source toolchains delivered via structured lab sequences improve accessibility and lower the barrier to entry for students without access to high-performance workstations \cite{b9}. The STM32 Nucleo platform has been validated as a pedagogically effective embedded systems teaching tool, offering substantial computational and peripheral advantages over traditional Arduino-class boards for IoT-oriented curricula \cite{b10}. On the technical side, Boolean reservoir computing implemented on FPGA achieves RF modulation classification accuracy competitive with convolutional neural networks (CNNs) while using a fraction of the trainable parameters \cite{b1}. Streaming CNN architectures deployed directly on RFSoC FPGAs have demonstrated continuous, sample-by-sample automatic modulation classification across eight modulation schemes at 128~MHz sampling frequency \cite{b3}. FPGA-accelerated angle-of-arrival estimation using ResNet reduced inference cycle count by 28.2\% compared to CPU execution while demonstrating robustness to multipath fading that degraded traditional MUSIC-based approaches \cite{b2}. Deep learning comparison studies have shown that DNNs can achieve 100\% modulation classification accuracy at signal-to-noise ratios of 4~dB and above on standard benchmark datasets \cite{b4}. For LoRa-specific signals, CNNs trained on spectrogram representations achieved 99.77\% spreading factor detection accuracy across a \(-\)20~dB to +30~dB SNR range \cite{b5}, while CNNs operating on raw I/Q data achieved 93.3\% accuracy for device identification under real experimental conditions \cite{b6}. Neural network-based LoRa demodulation from time-domain I/Q samples has been shown to match the performance of coherent matched-filter detection while tolerating carrier frequency offset and timing impairments that degrade conventional non-coherent receivers \cite{b8}. SoC FPGA implementations of LoRa gateways have achieved hardware-accelerated AES-128 encryption at 1200~Mbps while consuming only 17\% of available FPGA resources \cite{b7}. However, no existing curriculum integrates all four disciplines---FPGA design, RF communications, embedded systems, and machine learning---into a single progressive sequence in which every integration point is an explicit learning objective.

This paper addresses that gap. Section~\ref{sec:problem} formally defines the research problem and its context. Section~\ref{sec:framework} presents the proposed educational framework, including the hardware platform, curriculum module sequence, and pedagogical rationale.

\section{Problem Statement}
\label{sec:problem}

The research problem is the absence of a structured, progressive educational curriculum that integrates FPGA digital design, RF communications, embedded microcontroller programming, and machine learning-based signal analysis into a unified learning experience where each cross-domain interface is a deliberate teaching objective.

The target population is senior undergraduate and master's-level electrical engineering students who have completed introductory coursework in each of the four domains individually. These students possess the foundational prerequisites---HDL syntax, basic modulation and link theory, embedded C programming, and introductory machine learning---but lack experience working with systems that require all four simultaneously. When confronted with a multidisciplinary design problem, they face a secondary difficulty: understanding how a decision made within one domain generates requirements and constraints in adjacent domains that cannot be resolved from within a single discipline.

The concrete system context is a benchtop RF signal processing testbed. Multiple STM32 Nucleo-L476RG microcontrollers, each paired with an RFM95W LoRa transceiver, serve as independently configurable beacon transmitters. A receiver station captures incoming transmissions and forwards parsed RF data to a Digilent Nexys A7-100T FPGA, which performs real-time detection, packet integrity verification, and feature extraction in programmable logic on the Artix-7 fabric. Extracted signal features are collected by a host machine and used to train and evaluate machine learning classifiers for signal classification and emitter identification tasks. This architecture exposes four concrete interface problems that span disciplinary boundaries: the firmware-to-RF interface, the RF-to-FPGA interface, the FPGA-to-data-pipeline interface, and the data-to-model interface. Each requires cross-disciplinary reasoning to characterize and resolve.

The specific problem being solved is not a shortage of technical resources in any individual domain. Comprehensive textbooks, open-source FPGA IP cores, well-documented ML frameworks, and hardware reference designs exist for every component in this system. The gap is the absence of a curriculum that sequences these resources into an integrated learning experience, identifies the interface problems between domains as first-class subjects of instruction, and provides students with a concrete hardware implementation through which to develop and test their understanding.

\section{Proposed Integrated Educational Framework}
\label{sec:framework}

The proposed solution is a five-module curriculum paired with the benchtop hardware platform described in Section~\ref{sec:problem}. The framework assumes foundational domain knowledge and targets the integration layer: the design decisions, interface protocols, and system-level reasoning skills that emerge only when students work across disciplinary boundaries simultaneously.

\subsection{Hardware Platform}

The platform consists of three subsystem layers selected for their pedagogical properties alongside their technical capability.

At the transmission layer, up to four STM32 Nucleo-L476RG microcontrollers each drive an RFM95W LoRa transceiver module operating in the 915~MHz ISM band. The RFM95W is built on the Semtech SX1276 chipset and communicates with the host microcontroller over a standard 4-wire SPI bus, making it directly compatible with the STM32's hardware SPI peripheral. It supports spreading factors SF6 through SF12, signal bandwidths from 7.8~kHz to 500~kHz, and achieves a receiver sensitivity of $-$148~dBm at SF12 with 125~kHz bandwidth---parameters that students configure directly in firmware and observe as changes in link range, data rate, and noise tolerance. Each beacon transmits structured packets at configurable spreading factors, bandwidths, and power levels. The STM32 Nucleo-L476RG was selected in part for its demonstrated effectiveness as an embedded systems teaching platform: prior work reports that STM32 Nucleo-64 boards provide four independent SPI, I2C, and USART interfaces and include a built-in ST-LINK/V2-1 programmer, offering substantially greater peripheral capability than Arduino-class boards while remaining accessible to students learning embedded firmware \cite{b10}. The LoRa modulation scheme was chosen deliberately: its chirp spread spectrum waveform provides a set of directly observable RF parameters---spreading factor, bandwidth, coding rate---that connect modulation theory to both FPGA datapath design and machine learning feature engineering. The spreading factor in particular is a pedagogically rich parameter: increasing SF from 7 to 12 improves link budget by approximately 7.5~dB per step while halving the data rate, forcing students to reason about the tradeoff between link robustness and throughput in terms that directly affect downstream FPGA buffer sizing and ML dataset collection time. Experimental results confirm that CNNs trained on I/Q data from LoRa transmitters can reliably classify device identity and signal legitimacy from real hardware captures \cite{b6}, and that neural network-based demodulation of LoRa symbols from time-domain I/Q samples is robust to carrier frequency offset and timing impairments that challenge conventional receivers \cite{b8}.

At the processing layer, a receiver MCU forwards deserialized LoRa data to the Nexys A7-100T FPGA via UART or SPI through one of the board's twelve PMOD connectors. The Nexys A7-100T is built on a Xilinx Artix-7 XC7A100T device, providing 63,400 look-up tables (LUTs), 126,800 flip-flops, 240 DSP48E1 slices, and 4,860~Kb of block RAM across 135 tiles, along with six clock management tiles (CMTs) each containing a mixed-mode clock manager (MMCM) and a phase-locked loop (PLL). The board also provides 128~MB of DDR2 SDRAM and a 16~MB Quad-SPI flash for standalone operation. The FPGA performs packet parsing, CRC-based integrity checking, and feature extraction in RTL logic synthesized onto this fabric. The Artix-7 resource budget imposes realistic constraints that are tractable for a student implementing a first signal processing datapath: 240 DSP slices supports fixed-point arithmetic pipelines, while 4,860~Kb of block RAM is sufficient for packet FIFOs, feature buffers, and small on-chip weight storage. Prior work demonstrates that even computationally intensive inference pipelines fit comfortably within comparable logic budgets: a streaming CNN for modulation classification on an RFSoC FPGA consumed only approximately 10\% of available DSP slices and 15\% of block RAM while processing an uninterrupted 128~MHz data stream \cite{b3}, and a Boolean reservoir computing classifier matched CNN classification accuracy using far fewer trainable parameters \cite{b1}.

At the analysis layer, a host machine collects labeled feature vectors from the FPGA, builds a dataset across multiple beacon configurations, and trains ML classifiers. The choice of classification architecture is guided by comparative studies showing that CNNs and DNNs both achieve high accuracy on benchmark RF modulation datasets---with DNNs reaching 100\% accuracy for PAM, QAM, and PSK at SNR $\geq 4$~dB \cite{b4}---while LoRa-specific classifiers built on deep CNN architectures achieve 99.77\% spreading factor detection accuracy even at $-20$~dB SNR \cite{b5}. These results inform the model selection guidance provided to students in Module 4.

\subsection{Engagement Simulator}

To bridge the gap between curriculum introduction and physical hardware bring-up, a custom software engagement simulator was developed in Python as a standalone pre-hardware learning tool. The simulator models a multi-asset electronic warfare (EW) scenario---two cooperative target stations and an adversarial jamming threat---using a tick-based event engine that processes communication, jamming, and counter-jamming interactions at a configurable rate (500~ms per tick by default).

The simulator incorporates a physics-grounded RF model directly connected to the theoretical content students encounter in Module 1. Free-space path loss (FSPL) is computed per the standard Friis equation, and maximum communication range is derived from a link budget using configurable transmit power, antenna gain, and receiver sensitivity parameters. Target assets are modeled as frequency-hopping spread spectrum (FHSS) radios operating at 225~MHz (UHF military band) with eight available channels and a seed-synchronized hop sequence that both target stations share; the jamming threat uses a wideband receiver and DRFM-style noise jammer. This FHSS model connects directly to the spectral and synchronization concepts students work with during FPGA RTL design in Module 2, where clock domain management and frequency-selective filtering are primary design challenges.

Asset positions are specified in WGS-84 geodetic coordinates (latitude, longitude, altitude), with haversine-based range computation and geographic visualization on a 2-D map canvas. The default scenario places assets at coordinates representative of a realistic operational area, giving students immediate exposure to the link budget reasoning---how transmitted power, range, and receiver sensitivity interact---that governs the LoRa beacon configurations they will set in Module 1.

At each simulation tick, the engine resolves target-to-target communication attempts, threat jamming actions, and target counter-jamming responses in sequence. The outcome of each event is determined by the current channel alignment, effective range, and stochastic counter-jamming probability (default: 40\%). All events are written to a dual-log output: \texttt{engagement.log} records timestamped event descriptions, and \texttt{data\_stream.log} records the raw binary sequences exchanged---structured DATA, ACK, JAM, and counter-JAM patterns. These labeled logs serve as the first training dataset available to students in Module 4, allowing ML pipeline development to begin before any hardware RF data has been collected. Critically, the ground-truth labels (communication success, jamming state, counter-jamming outcome) are embedded directly in the log structure, giving students a clean supervised learning problem on which to validate preprocessing and model training code before the noisier, real-world dataset from the FPGA is available.

Per-tick performance metrics---communication success rate, jamming effectiveness, and counter-jamming rate---are accumulated over a rolling 60-tick history and displayed as live charts. These same three metrics form the primary benchmarks students measure on the physical system in Module 5, so the simulator establishes the measurement methodology and expected performance envelope before students have seen any hardware results.

\subsection{Curriculum Module Sequence}

\textit{Module 1: RF Hardware and Subsystem Design.} Students implement beacon transmission firmware on the STM32 platform, design the packet protocol and transmission scheduling algorithm, and validate the physical receiver-to-FPGA serial interface. The primary learning objective is understanding how firmware design choices---SPI clock rate, DMA configuration, interrupt latency---affect the timing and integrity of the data stream delivered to the FPGA. The Nucleo-L476RG's four independent serial interfaces and hardware-accelerated peripherals \cite{b10} provide enough configuration space to make these tradeoffs concrete and experimentally observable.

\textit{Module 2: FPGA RTL Design and Simulation.} Students implement the FPGA-side data ingestion pipeline in VHDL or Verilog: serial receiver, packet parser, CRC checker, and a feature extraction datapath. Simulation testbenches use recorded beacon data from Module 1 as stimulus. Students encounter FPGA-specific design constraints---timing closure, clock domain crossing, memory resource allocation---in the context of a real signal protocol rather than a synthetic exercise. Cloud-based FPGA toolchains have demonstrated that structured, hardware-grounded lab sequences can make this learning accessible without requiring specialized hardware beyond a web browser \cite{b9}; this module follows that model while adding the RF protocol context that a standalone digital design course omits.

\textit{Module 3: System Integration and Data Collection.} Students connect the hardware subsystems built in Modules 1 and 2 and characterize end-to-end data flow: from RF beacon transmission through reception and FPGA processing to feature output. Latency, throughput, and packet error rate are measured. Synchronization mismatches and protocol errors that were invisible in per-subsystem development surface here as integration failures that must be diagnosed and resolved. Students then execute a structured data collection campaign across multiple beacon configurations, producing a labeled RF dataset. This module deliberately surfaces the distribution shift problems that arise when training data is drawn from real hardware rather than simulation \cite{b6}.

\textit{Module 4: AI-Driven Classification Pipeline.} Using the hardware-collected dataset from Module 3, students preprocess FPGA-extracted features, train and evaluate classifiers, and analyze classification accuracy as a function of SNR and number of concurrent beacons. Streaming inference architectures \cite{b3} and FPGA-accelerated ML pipelines \cite{b2} serve as reference designs for students considering on-device deployment. Comparative architecture studies \cite{b4} and LoRa-specific classification results \cite{b5,b6,b8} provide empirical benchmarks against which students evaluate their own model performance.

\textit{Module 5: System Validation and Benchmarking.} Students stress-test the complete pipeline with multiple simultaneous beacons, measure total system response time from RF transmission to classification decision, and characterize the performance envelope of their implementation. The benchmarking exercise requires active engagement across all four domains: RF link budget analysis, FPGA timing and resource reports, serial interface throughput, and ML classification accuracy curves.

\subsection{Pedagogical Rationale}

The framework's central design principle is that the interfaces between disciplinary domains are the primary teaching objects. Existing FPGA education frameworks \cite{b9} and embedded systems curricula \cite{b10} have demonstrated the effectiveness of hardware-grounded, progressive lab sequences within a single discipline. This work extends that model to four disciplines simultaneously by treating each inter-domain interface as a design problem whose resolution requires knowledge from both adjacent domains.

A student who implements a CRC verifier in RTL and then diagnoses a packet loss event caused by firmware DMA timing has internalized a cross-domain constraint that no single-discipline course can produce. Likewise, a student who trains a classifier on FPGA-extracted features and observes accuracy degradation when beacon power levels change has encountered the coupling between RF link behavior and machine learning model performance as an empirical reality rather than a textbook abstraction.

All hardware components are available through standard academic purchasing channels, and all development tools are free or available under academic license. Prior work has shown that hardware-grounded RF and signal classification pipelines are feasible within the resource budgets of mid-range educational FPGAs \cite{b1,b3,b7}, establishing the technical viability of the platform. The framework's contribution is not a novel hardware design or machine learning technique but a documented, teachable model of how FPGA design, RF communications, embedded firmware, and machine learning interact at the system level---organized as a progressive curriculum in which that interaction is the subject of instruction.

\begin{thebibliography}{00}

\bibitem{b1}
H. Komkov, L. Pocher, A. Restelli, B. Hunt, and D. Lathrop, ``RF signal classification using Boolean reservoir computing on an FPGA,'' in \textit{Proc. 2021 Int. Joint Conf. Neural Networks (IJCNN)}, Jul. 2021, doi: 10.1109/IJCNN52387.2021.9533342.

\bibitem{b2}
J. F. Lin, T. Morehouse, C. Montes, E. Caushi, A. Dudko, E. Savage, and R. Zhou, ``A study of angle of arrival estimation of a RF signal with FPGA acceleration,'' in \textit{Proc. 2024 Int. Wireless Commun. Mobile Comput. Conf. (IWCMC)}, 2024, pp. 1625--1628, doi: 10.1109/IWCMC61514.2024.10592551.

\bibitem{b3}
A. Maclellan, L. H. Crockett, and R. W. Stewart, ``Streaming convolutional neural network FPGA architecture for RFSoC data converters,'' in \textit{Proc. 2023 21st IEEE Interreg. NEWCAS Conf. (NEWCAS)}, Edinburgh, UK, Jun. 2023, doi: 10.1109/NEWCAS57931.2023.10198198.

\bibitem{b4}
K. A. Alnajjar, S. Ghunaim, and S. Ansari, ``Automatic modulation classification in deep learning,'' in \textit{Proc. 2022 5th Int. Conf. Commun., Signal Process., Appl. (ICCSPA)}, Cairo, Egypt, Dec. 2022, doi: 10.1109/ICCSPA55860.2022.10019122.

\bibitem{b5}
P. M. Mutescu, A. Lavric, A. I. Petrariu, and V. Popa, ``Deep learning enhanced spectrum sensing for LoRa spreading factor detection,'' in \textit{Proc. 2023 13th Int. Symp. Adv. Topics Elect. Eng. (ATEE)}, Bucharest, Romania, Mar. 2023, doi: 10.1109/ATEE58038.2023.10108224.

\bibitem{b6}
Y. E. Sagduyu and T. Erpek, ``Adversarial attacks on LoRa device identification and rogue signal detection with deep learning,'' in \textit{Proc. MILCOM 2023 --- 2023 IEEE Military Commun. Conf.}, Boston, MA, Oct.--Nov. 2023, pp. 385--390, doi: 10.1109/MILCOM58377.2023.10356224.

\bibitem{b7}
T.-P. Dang, T.-K. Tran, T.-T. Bui, and H.-T. Huynh, ``LoRa gateway based on SoC FPGA platforms,'' in \textit{Proc. 2021 Int. Symp. Elect. Electron. Eng. (ISEE)}, 2021, pp. 48--52, doi: 10.1109/ISEE51682.2021.9418711.

\bibitem{b8}
K. Dakic, B. Al Homssi, A. Al-Hourani, and M. Lech, ``LoRa signal demodulation using deep learning, a time-domain approach,'' in \textit{Proc. 2021 IEEE 93rd Veh. Technol. Conf. (VTC2021-Spring)}, Apr. 2021, doi: 10.1109/VTC2021-Spring51267.2021.9448711.

\bibitem{b9}
W. Smith, Z. Driskill, J. Goeders, and M. Wirthlin, ``Digital design education using an open-source, cloud-based FPGA toolchain,'' in \textit{Proc. 2024 Intermountain Eng., Technol. Comput. Conf. (IETC)}, May 2024, doi: 10.1109/IETC61393.2024.10564285.

\bibitem{b10}
N. O. Strelkov, V. V. Krutskikh, and E. V. Shalimova, ``Programming STM32 Nucleo platform for IoT education using STM32duino and Mbed OS,'' in \textit{Proc. 2022 VI Int. Conf. Inf. Technol. Eng. Educ. (Inforino)}, Apr. 2022, doi: 10.1109/Inforino53888.2022.9782920.

\end{thebibliography}

\end{document}
